\section{The agent-usability framework}
\label{sec:framework}

\subsection{Design rationale}

An agent is not a data consumer in the classical sense. It does not integrate
two systems once, at engineering time, with a human present to resolve
ambiguity. It receives a natural-language goal, must resolve it against a plant
it has never seen, must decide what evidence it needs, must act, and must be
prevented from acting wrongly, continuously, unsupervised, at inference time.
The framework below asks, of each technology, what it contributes to that loop.

Each criterion is stated as a capability the technology either grants or
withholds, together with the concrete agent failure it prevents. The four-point
rubric is in Appendix~\ref{app:rubric}.

\subsection{The ten criteria}

\paragraph{U1: Identity and resolvability.}
Does the technology assign globally unique, stable, ideally dereferenceable
identifiers to concepts \emph{and} to instances, and can an agent get from a
human phrase to such an identifier? IRDIs and IRIs do this; a historian tag name
does not, because it is unique only inside one system and its meaning is carried
by an unwritten naming convention. \emph{Prevents:} referring to the wrong
physical thing; silently conflating homographs.

\paragraph{U2: Type system and schema exposure.}
Are classes, properties, datatypes, cardinalities, value ranges, enumerations
and access rights machine-readable, and can the agent \emph{introspect} them at
run time rather than being told about them out of band? \emph{Prevents:}
malformed calls; writes to read-only variables; out-of-range setpoints;
invented properties.

\paragraph{U3: Quantity and unit semantics.}
Is a value bound to a \emph{quantity kind} and a \emph{unit} with a declared
dimension and conversion, or merely to a display string? A unit symbol is
decoration; a dimension vector plus an affine conversion is a calculus.
\emph{Prevents:} the Mars-Climate-Orbiter class of error, meaning dimensionally
absurd assignments, and worse, dimensionally \emph{correct} assignments at the
wrong scale, which no type checker catches.

\paragraph{U4: Relational expressivity.}
Can the model express mereology (part--whole, ISA-95 hierarchy), topology
(material and utility flow), and arbitrary typed relations, with transitivity
where it applies? \emph{Prevents:} inability to scope a query
(``\emph{Line~1} tags''), to propagate impact (``if the compressor trips''), or
to authorise by containment.

\paragraph{U5: Behavioural and action semantics.}
Does the model describe what may be \emph{done}: callable operations with
signatures, state machines, service interfaces, pre- and post-conditions?
\emph{Prevents:} commands that are illegal in the current state; plans whose
steps are individually plausible and jointly impossible.

\paragraph{U6: Constraint and validation semantics.}
Can the technology express constraints in a form an engine can \emph{check}, as
opposed to prose in a manual? SHACL shapes, OWL disjointness, interlock guards
and MTP conditions qualify; a PDF specification does not. \emph{Prevents:}
everything else, after the fact. This is the criterion that converts an
LLM's output from a claim into a proposition with a truth value.

\paragraph{U7: Provenance, time and trust.}
Are observations timestamped, attributed to a source, and versioned; is
uncertainty or quality representable? \emph{Prevents:} decisions on stale data;
unfalsifiable answers; inability to audit an agent's reasoning after an
incident.

\paragraph{U8: Queryability and aggregation pushdown.}
Is there a standard query language, and can filtering, joining, aggregation and
transitive traversal be evaluated \emph{outside} the model, returning only the
answer? \emph{Prevents:} the fundamental scaling failure of context stuffing
(Experiment~E7).

\paragraph{U9: Context economy.}
What does the representation cost in tokens per unit of usable fact, and does
it support projection, so that a consumer can request a small, faithful slice?
\emph{Prevents:}
budget exhaustion before the relevant evidence is reached, the dominant cause of
apparent ``reasoning'' failures in long-context agents.

\paragraph{U10: Ecosystem and deployability.}
Are there open implementations, converters, public reference data, and a
governance body? A technology that scores well on U1--U9 and has no runtime is a
research proposal. \emph{Prevents:} the project never happening.

\subsection{Non-substitutability}

The framework's substantive claim is that U1--U8 are \emph{orthogonal}: strength
in one does not compensate for weakness in another. This is not obvious. One
might reasonably argue that a sufficiently rich type system implies unit safety
(it does not: \texttt{Double} in \texttt{bar} and \texttt{Double} in \texttt{psi}
are the same type), or that a complete information model implies behavioural
safety (it does not: an \opcua{} address space describes the variables of a
machine, not the order in which its commands are legal).

We test the claim directly. The five context conditions of \S\ref{sec:design}
add criteria cumulatively, and the five validator tiers do the same. If the
criteria were substitutable, the curves would saturate early. Instead
(Figure~\ref{fig:ceiling}, Figure~\ref{fig:guardrails}) each added criterion
opens a new band of tasks and a new family of detectable faults, and the
budget sweep shows that no amount of the criteria a representation \emph{has}
substitutes for one it lacks.

\subsection{The comparison matrix}

Table~\ref{tab:auf-matrix} applies the rubric across the survey. Reading it by
column rather than by row is instructive.

\begin{table*}[t]
\centering
\caption{Agent-usability assessment. Scores $0$--$3$ per the rubric in
Appendix~\ref{app:rubric}, assessed against what the specification standardises,
not what individual tools add. \textbf{U1} identity, \textbf{U2} types,
\textbf{U3} units, \textbf{U4} relations, \textbf{U5} behaviour,
\textbf{U6} constraints, \textbf{U7} provenance, \textbf{U8} query,
\textbf{U9} economy, \textbf{U10} ecosystem.}
\label{tab:auf-matrix}
\footnotesize
\setlength{\tabcolsep}{4pt}
\resizebox{\ifdim\width>\linewidth\linewidth\else\width\fi}{!}{%
\begin{tabular}{@{}llcccccccccc@{}}
\toprule
\textbf{L} & \textbf{Technology} & U1 & U2 & U3 & U4 & U5 & U6 & U7 & U8 & U9 & U10 \\
\midrule
L0 & ECLASS / IEC CDD / IEC 61987      & 3 & 2 & 2 & 1 & 0 & 1 & 0 & 1 & 2 & 3 \\
L0 & GS1 (GTIN/GIAI/EPCIS)             & 3 & 1 & 0 & 1 & 1 & 0 & 2 & 1 & 2 & 3 \\
\midrule
L1 & \opcua{} core + companions (IEC 62541) & 3 & 3 & 2 & 2 & 3 & 1 & 2 & 1 & 1 & 3 \\
L1 & Asset Administration Shell (IDTA) & 3 & 3 & 2 & 2 & 2 & 1 & 2 & 1 & 1 & 3 \\
L1 & AutomationML (IEC 62714)          & 2 & 3 & 1 & 3 & 1 & 1 & 1 & 0 & 1 & 2 \\
L1 & MTP (VDI/VDE/NAMUR 2658)          & 2 & 3 & 2 & 2 & 3 & 2 & 1 & 0 & 2 & 2 \\
L1 & PackML (ISA-TR88.00.02)           & 1 & 2 & 0 & 0 & 3 & 2 & 0 & 0 & 3 & 3 \\
L1 & ISA-95 / IEC 62264 (B2MML)        & 2 & 2 & 1 & 3 & 1 & 0 & 1 & 0 & 2 & 3 \\
L1 & W3C WoT Thing Description 1.1     & 3 & 3 & 2 & 1 & 3 & 1 & 1 & 1 & 2 & 2 \\
L1 & NGSI-LD (ETSI GS CIM 009)         & 3 & 2 & 2 & 3 & 1 & 1 & 3 & 2 & 2 & 2 \\
L1 & IEC 61850                         & 3 & 3 & 2 & 2 & 2 & 1 & 2 & 0 & 1 & 3 \\
\midrule
L2 & RDF / RDFS                        & 3 & 2 & 0 & 3 & 0 & 0 & 1 & 3 & 1 & 3 \\
L2 & OWL 2 (EL/QL/RL/DL)               & 3 & 3 & 0 & 3 & 0 & 2 & 0 & 2 & 1 & 3 \\
L2 & SHACL                             & 2 & 3 & 1 & 2 & 1 & 3 & 1 & 2 & 2 & 3 \\
L2 & SPARQL 1.1                        & 2 & 1 & 0 & 3 & 0 & 1 & 1 & 3 & 3 & 3 \\
L2 & PROV-O                            & 3 & 2 & 0 & 3 & 1 & 0 & 3 & 2 & 2 & 2 \\
\midrule
L3 & QUDT                              & 3 & 2 & 3 & 1 & 0 & 2 & 0 & 2 & 2 & 2 \\
L3 & SOSA / SSN                        & 3 & 2 & 1 & 2 & 1 & 0 & 3 & 2 & 2 & 3 \\
L3 & Brick                             & 3 & 3 & 1 & 3 & 0 & 1 & 0 & 3 & 2 & 3 \\
L3 & SAREF (+ extensions)              & 3 & 2 & 2 & 2 & 2 & 0 & 1 & 2 & 2 & 2 \\
L3 & DEXPI                             & 2 & 3 & 2 & 3 & 0 & 1 & 1 & 1 & 1 & 2 \\
\midrule
L4 & IDS-RAM / Dataspace Protocol      & 3 & 2 & 0 & 2 & 1 & 2 & 3 & 1 & 1 & 2 \\
L4 & Catena-X / Manufacturing-X        & 3 & 3 & 1 & 2 & 1 & 2 & 3 & 1 & 1 & 2 \\
L4 & ODRL                              & 2 & 2 & 0 & 2 & 2 & 3 & 2 & 1 & 2 & 2 \\
\bottomrule
\end{tabular}}
\end{table*}

Three column-wise observations drive the rest of the paper.

\begin{itemize}
  \item \textbf{U2 (types) is solved.} Eleven of the surveyed technologies score
  $3$. This is the industry's genuine achievement and it is why so many
  ``AI-ready plant data'' offerings appear convincing: they surface real types.
  \item \textbf{U3 (units), U5 (behaviour) and U6 (constraints) are sparse.}
  Only QUDT scores $3$ on units; only \opcua{} methods, WoT actions, MTP and
  PackML score $3$ on behaviour; only SHACL and ODRL score $3$ on constraints.
  Crucially, \emph{no single technology scores $3$ on all three}, so an agent
  that needs all three, and \S\ref{sec:results} shows it does, must
  compose several standards.
  \item \textbf{U9 (economy) is inversely correlated with U2--U7.} The richest
  models are the most verbose. This is not a criticism of the standards; it is a
  statement that the semantic layer belongs behind a query interface, not in the
  prompt.
\end{itemize}
